\documentclass[12pt,a4paper]{article}

\usepackage[utf8]{inputenc}
\usepackage[T1]{fontenc}
\usepackage[brazil]{babel}
\usepackage{amsmath,amssymb,amsfonts}
\usepackage{geometry}
\geometry{top=3cm,bottom=2cm,left=3cm,right=2cm}
\usepackage{microtype}
\usepackage{booktabs}
\usepackage{graphicx}
\usepackage{setspace}
\onehalfspacing
\usepackage{indentfirst}
\usepackage{hyperref}
\hypersetup{
    colorlinks=true,
    linkcolor=black,
    citecolor=black,
    urlcolor=blue
}
\usepackage[ruled,vlined,linesnumbered,portuguese]{algorithm2e}
\usepackage{cite}

% Marcador temporário para trechos que dependem de dados, implementação ou decisão.
% Remover todas as ocorrências antes da versão final.
\newcommand{\pendente}[1]{\textbf{[PENDENTE: #1]}}

\begin{document}

% --- ELEMENTOS PRÉ-TEXTUAIS ---

\begin{titlepage}
\begin{center}
    {\large \textbf{UNIVERSIDADE FEDERAL FLUMINENSE}}\\[0.2cm]
    {\large \textbf{INSTITUTO DE COMPUTAÇÃO}}\\[0.2cm]
    {\large \textbf{DEPARTAMENTO DE CIÊNCIA DA COMPUTAÇÃO}}\\[4cm]

    {\large \textbf{NOME DO ALUNO}}\\[3.5cm]

    {\Large \textbf{Metaheurísticas Aplicadas à Alocação de Horários e Salas de Aula no Instituto de Computação da UFF}}\\[4cm]

    \vfill

    {\large Niterói}\\[0.2cm]
    {\large 2026}
\end{center}
\end{titlepage}

\newpage

\begin{center}
    {\large \textbf{NOME DO ALUNO}}\\[3cm]
    {\Large \textbf{Metaheurísticas Aplicadas à Alocação de Horários e Salas de Aula no Instituto de Computação da UFF}}\\[3cm]
\end{center}

\vspace{2cm}
\hfill
\begin{minipage}{0.5\textwidth}
    Trabalho de Conclusão de Curso submetido ao Departamento de Ciência da Computação da Universidade Federal Fluminense como requisito parcial para a obtenção do título de Bacharel em Ciência da Computação.\\[0.5cm]
    \textbf{Orientador:} Prof. Dr. Igor Machado Coelho
\end{minipage}

\vfill

\begin{center}
    {\large Niterói}\\[0.2cm]
    {\large 2026}
\end{center}

\newpage

% --- RESUMO E ABSTRACT ---

\begin{center}
    \textbf{\large Resumo}
\end{center}
\vspace{0.5cm}

A elaboração de quadros de horários universitários envolve decisões combinatórias sobre turmas, professores, períodos e salas, sujeitas a restrições de viabilidade e critérios de qualidade potencialmente conflitantes. Este trabalho aborda o problema de atribuição de professores, horários e salas no Instituto de Computação da Universidade Federal Fluminense (IC/UFF). O modelo considera um horizonte anual, as grades curriculares de Ciência da Computação e Sistemas de Informação, turmas compartilhadas e a alocação de sala por encontro. A abordagem computacional prevê componentes de representação, avaliação e vizinhança e o emprego de metaheurísticas por meio do \textit{pyoptframe}. Uma pipeline de dados foi construída a partir dos quadros de horários e das grades curriculares. \pendente{completar o resumo com as instâncias definitivas, os algoritmos comparados e os principais resultados experimentais.}

\vspace{0.5cm}
\noindent \textbf{Palavras-chave}: Timetabling Universitário, Metaheurísticas, OptFrame, Alocação de Salas, Otimização Combinatória.

\vspace{1.5cm}

\begin{center}
    \textbf{\large Abstract}
\end{center}
\vspace{0.5cm}

University course timetabling involves combinatorial decisions concerning classes, instructors, periods, and rooms, subject to feasibility constraints and potentially conflicting quality criteria. This work addresses the assignment of instructors, times, and rooms at the Institute of Computing of Fluminense Federal University (IC/UFF). The model considers an annual planning horizon, the Computer Science and Information Systems curricula, shared classes, and room assignment per class meeting. The computational approach is designed around representation, evaluation, and neighborhood components and the use of metaheuristics through \textit{pyoptframe}. A data pipeline was built from historical timetables and curriculum data. \textbf{[PENDING: complete the abstract with the final instances, compared algorithms, and main experimental results.]}

\vspace{0.5cm}
\noindent \textbf{Keywords}: University Course Timetabling, Metaheuristics, OptFrame, Room Allocation, Combinatorial Optimization.

\newpage
\tableofcontents
\newpage

% --- CORPO DO TRABALHO ---

\section{Introdução}
\label{sec:introducao}

A elaboração de quadros de horários e a alocação de salas em instituições de ensino superior constituem problemas de otimização combinatória estudados na literatura de \textit{educational timetabling}. No caso universitário, esse conjunto de problemas é usualmente denominado \textit{University Course Timetabling Problem} (UCTP) \cite{schaerf1999survey,lewis2008survey}. Mesmo formulações simplificadas contêm problemas NP-difíceis como casos particulares, além de combinarem restrições de viabilidade e critérios de qualidade distintos.

No Instituto de Computação da Universidade Federal Fluminense (IC/UFF), a modelagem considerada neste trabalho envolve três papéis institucionais. A chefia de departamento e a coordenação participam da distribuição de encargos e da organização dos horários; a coordenação também representa a necessidade dos alunos de cursar as disciplinas previstas sem choques; e o Instituto administra as salas e recebe demandas de capacidade e recursos. Esses papéis compartilham recursos, mas avaliam a solução sob perspectivas diferentes.

O problema considerado inclui padrões de horário por semestre par e ímpar, agrupamento de disciplinas em setores associados a dias da semana, horários fixos para disciplinas prioritárias ou externas e critérios referentes a preferências docentes, quantidade de dias com aula, janelas e descanso entre jornadas. Também são consideradas uma carga anual mínima de três disciplinas obrigatórias por professor e as grades de Ciência da Computação (CC) e Sistemas de Informação (SI), inclusive as turmas compartilhadas. \pendente{inserir a fonte institucional das regras de carga anual, prioridade e descanso, além de delimitar o universo de docentes abrangidos.}

A combinação dessas decisões motiva o uso de métodos computacionais de apoio à elaboração e à avaliação dos quadros de horários. O propósito do trabalho é construir e avaliar um modelo próximo da realidade observada, documentando as hipóteses adotadas e a origem dos dados utilizados.

\subsection{Objetivos}
\label{subsec:objetivos}

O objetivo geral deste trabalho é modelar o problema integrado de atribuição de professores, horários e salas de aula no IC/UFF e investigar o uso de metaheurísticas por meio do \textit{pyoptframe} para a geração e a avaliação de soluções.

Os objetivos específicos são:
\begin{enumerate}
    \item formalizar o problema com base no \textit{Curriculum-Based Course Timetabling} (CB-CTT), adaptando-o à atribuição de professores, ao horizonte anual e às regras do IC/UFF;
    \item estruturar uma pipeline de dados a partir dos quadros de horários, das grades curriculares e de outros dados institucionais disponíveis;
    \item implementar a representação da solução, o avaliador e os movimentos de sala, horário e professor;
    \item implementar e comparar as metaheurísticas selecionadas por meio do \textit{pyoptframe};
    \item executar um protocolo com múltiplas sementes e comparar as soluções com um baseline histórico avaliado pelos mesmos critérios;
    \item estudar o efeito da ordem de planejamento das grades de CC e SI por meio dos cenários E1, E2 e E3 descritos na Seção~\ref{sec:experimentos}.
\end{enumerate}

\subsection{Organização do Trabalho}
\label{subsec:organizacao}

O restante deste texto está organizado da seguinte forma. A Seção~\ref{sec:fundamentacao} apresenta a fundamentação sobre timetabling, metaheurísticas e OptFrame. A Seção~\ref{sec:modelagem} formaliza o problema. A Seção~\ref{sec:proposta_solucao} descreve a arquitetura computacional, a representação, o avaliador e os movimentos. A Seção~\ref{sec:experimentos} apresenta as instâncias, o protocolo experimental, os cenários de comparação e os resultados. Por fim, a Seção~\ref{sec:conclusao} reúne as conclusões e possibilidades de trabalhos futuros.

\section{Fundamentação Teórica}
\label{sec:fundamentacao}

Esta seção apresenta os conceitos necessários para situar a proposta: problemas de timetabling educacional, a formulação CB-CTT, metaheurísticas de trajetória e o framework OptFrame.

\subsection{Problemas de Timetabling Educacional}
\label{subsec:timetabling_teoria}

O termo \textit{timetabling} designa a atribuição de eventos a períodos e, conforme a formulação, a outros recursos, sujeita a restrições de viabilidade e critérios de qualidade. Schaerf \cite{schaerf1999survey} organiza os problemas educacionais em três categorias principais:
\begin{enumerate}
    \item \textbf{\textit{School Timetabling}}: atribuição periódica de aulas a professores e turmas escolares;
    \item \textbf{\textit{Exam Timetabling}}: agendamento de exames, com conflitos definidos principalmente pelos estudantes que realizam as provas;
    \item \textbf{\textit{University Course Timetabling}}: agendamento periódico de aulas universitárias, considerando professores, grupos de estudantes, períodos e salas.
\end{enumerate}

As definições e restrições exatas variam entre instituições. Por isso, UCTP designa uma família de problemas, e não uma única formulação universal \cite{lewis2008survey,babaei2015survey}.

\subsubsection{Curriculum-Based Course Timetabling}

Duas formulações recorrentes diferem pela origem dos conflitos discentes:
\begin{itemize}
    \item no \textit{Post-Enrolment Course Timetabling}, os conflitos são derivados das matrículas ou escolhas individuais dos estudantes;
    \item no \textit{Curriculum-Based Course Timetabling} (CB-CTT), os conflitos são definidos por \textit{curricula}, isto é, grupos de disciplinas cujos alunos se sobrepõem.
\end{itemize}

A trilha 3 da ITC-2007 consolidou uma formulação de referência para CB-CTT \cite{digaspero2007itc}. Nela, uma solução atribui a cada aula um período e uma sala. Entre as restrições fortes estão indisponibilidades, conflitos de professor ou currículo e ocupação exclusiva de sala. A capacidade da sala, o número mínimo de dias, a compactação curricular e a estabilidade de sala aparecem como critérios suaves. Uma visão posterior da literatura de CB-CTT é apresentada por Bettinelli et al. \cite{bettinelli2015overview}.

O presente trabalho usa o CB-CTT como referência, mas não o reproduz literalmente. Em particular, a atribuição do professor é uma decisão do modelo do IC/UFF, enquanto no benchmark da ITC-2007 o professor já pertence ao curso. Os grupos curriculares de CC e SI são construídos a partir das disciplinas que devem ser cursáveis em conjunto em cada período. \pendente{completar o mapeamento dos grupos curriculares e das turmas compartilhadas.}

\subsubsection{Complexidade Computacional}

Uma versão simplificada do problema de horários contém o problema de coloração de vértices. Considere um grafo de conflitos $G=(V,E)$, no qual cada vértice representa um evento e cada aresta indica que dois eventos não podem ocorrer simultaneamente. Associar um dos $K$ períodos a cada evento equivale a obter uma $K$-coloração:
\begin{equation}
    c(u) \neq c(v), \qquad \forall (u,v)\in E.
\end{equation}
Como decidir a existência de uma $K$-coloração é NP-completo para $K\geq 3$ \cite{garey1979computers}, essa versão de decisão pode ser reduzida ao problema de horários. Consequentemente, as formulações de otimização que a generalizam são NP-difíceis. Restrições adicionais de salas, recursos e preferências não eliminam essa dificuldade.

\subsection{Metaheurísticas de Otimização Combinatória}
\label{subsec:metaheuristicas_teoria}

Métodos exatos podem ser apropriados para determinadas formulações e tamanhos de instância, mas podem enfrentar dificuldades de escala em problemas combinatórios grandes. Metaheurísticas oferecem outra estratégia: explorar o espaço de soluções sem garantir otimalidade, buscando soluções de boa qualidade dentro do orçamento computacional disponível \cite{talbi2009metaheuristics}.

Seja $\mathcal{S}$ o espaço de soluções e $F:\mathcal{S}\rightarrow\mathbb{R}$ uma função de avaliação a minimizar. Uma estrutura de vizinhança $N$ associa a cada solução $s$ um conjunto de soluções alcançáveis por movimentos elementares.

\subsubsection{Simulated Annealing}

O \textit{Simulated Annealing} (SA), proposto por Kirkpatrick, Gelatt e Vecchi \cite{kirkpatrick1983optimization}, admite ocasionalmente movimentos de piora para reduzir a chance de aprisionamento precoce em ótimos locais. Para uma solução atual $s$ e uma vizinha $s'$, define-se:
\begin{equation}
    \Delta = F(s')-F(s).
\end{equation}
Movimentos com $\Delta\leq 0$ são aceitos. Quando $\Delta>0$, a aceitação ocorre com probabilidade:
\begin{equation}
    P(\text{aceitação})=\exp\left(-\frac{\Delta}{T}\right),
\end{equation}
onde $T>0$ é a temperatura. A temperatura é reduzida ao longo da execução. Esquemas geométricos, lineares e adaptativos são possibilidades distintas; a escolha e seus parâmetros devem ser declarados no protocolo experimental. À medida que a temperatura diminui, também diminui a probabilidade de aceitar uma piora de magnitude fixa.

\subsubsection{Iterated Local Search e Variable Neighborhood Search}

O \textit{Iterated Local Search} (ILS) alterna busca local, perturbação e um critério de aceitação para explorar diferentes ótimos locais \cite{lourenco2003iterated}. O \textit{Variable Neighborhood Search} (VNS) explora sistematicamente mais de uma estrutura de vizinhança, usando mudanças de vizinhança para combinar intensificação e diversificação \cite{hansen2001variable}.

\subsection{OptFrame e pyoptframe}
\label{subsec:optframe_teoria}

O OptFrame é um framework para o desenvolvimento de métodos de otimização combinatória com componentes reutilizáveis \cite{coelho2010optframe,coelho2020microbenchmark}. Entre esses componentes estão a representação da solução, o método construtivo, o avaliador, os movimentos, as estruturas de vizinhança e os métodos de busca. O \textit{pyoptframe} fornece bindings Python para o núcleo do OptFrame \cite{pyoptframe2026}.

A separação entre componentes permite alterar uma metaheurística ou uma vizinhança sem incorporar ao algoritmo genérico todos os detalhes do problema. Essa propriedade motiva o uso do framework neste trabalho e permite que representação, avaliação e movimentos sejam testados separadamente antes de sua composição nos métodos de busca.

\subsection{Regras institucionais consideradas}
\label{subsec:diretrizes_institucionais}

A modelagem considera regras locais de carga anual mínima de obrigatórias, prioridade entre professores, horários associados a setores e descanso entre jornadas. Essas regras complementam as restrições usuais de timetabling e afetam tanto os domínios das variáveis quanto a avaliação das soluções.

\pendente{inserir a referência normativa aplicável e definir formalmente: o universo de docentes da carga anual; a unidade de contagem da carga mínima; o critério de prioridade; e o intervalo mínimo de descanso.}

\section{Modelagem Matemática do Problema}
\label{sec:modelagem}

Esta seção formaliza o problema de atribuição de professores, horários e salas considerado no trabalho. As decisões ainda não fechadas são indicadas por marcadores explícitos para posterior preenchimento.

\subsection{Premissas e decisões de modelagem}
\label{subsec:premissas_modelagem}

O problema é tratado como uma variante de CB-CTT. A formulação de referência já inclui a atribuição de períodos e salas; as principais adaptações consideradas aqui são a atribuição de professores, as regras de setores e horários fixos, a alocação de sala por encontro e o horizonte anual.

O horizonte reúne os semestres ímpar e par porque a carga anual de obrigatórias e o rodízio relacionam decisões dos dois semestres. Restrições de conflito por horário são verificadas separadamente em cada semestre.

O modelo considera as grades de CC e SI. Professores e salas são recursos compartilhados. Uma disciplina comum às duas grades pode corresponder a uma turma única compartilhada. \pendente{inserir o mapeamento definitivo das turmas compartilhadas entre CC e SI.}

As turmas são classificadas por dois eixos:
\begin{enumerate}
    \item \textbf{responsabilidade}: turmas do IC ($\mathcal{C}^{\mathrm{IC}}$), cujas atribuições são decisões do modelo, e turmas de outros departamentos ($\mathcal{C}^{\mathrm{out}}$), tratadas como ocupações exógenas;
    \item \textbf{papel curricular}: obrigatórias ($\mathcal{C}^{\mathrm{obr}}$) e optativas ($\mathcal{C}^{\mathrm{opt}}$).
\end{enumerate}

O subconjunto $\mathcal{C}^{\mathrm{fix}}$ reúne turmas com horário dado, incluindo as externas e as disciplinas do IC que tenham prioridade ou padrão fixo confirmado.

\subsection{Conjuntos e parâmetros}
\label{subsec:conjuntos_parametros}

\begin{itemize}
    \item $\mathcal{C}$: turmas ofertadas no horizonte anual, indexadas por $c$;
    \item $\mathcal{P}$: professores do IC, indexados por $p$;
    \item $\mathcal{P}^{\mathrm{H12}}\subseteq\mathcal{P}$: docentes sujeitos à carga anual H12; \pendente{definir a composição de $\mathcal{P}^{\mathrm{H12}}$;}
    \item $\mathcal{R}$: salas consideradas, indexadas por $r$;
    \item $\mathcal{D}$: dias letivos da semana;
    \item $\mathcal{H}$: slots de horário, indexados por $h$;
    \item $\mathcal{Q}$: padrões semanais de horário, indexados por $q$, com slots $H_q\subseteq\mathcal{H}$;
    \item $\mathcal{S}$: setores do departamento;
    \item $\mathcal{K}=\{\mathrm{CC},\mathrm{SI}\}$: grades curriculares;
    \item $\mathcal{G}$: grupos curriculares; cada grupo está associado a uma grade e a um semestre;
    \item $\mathcal{M}_c$: encontros semanais da turma $c$.
\end{itemize}

Mapeamentos e domínios:
\begin{itemize}
    \item $\sigma(c)\in\{\text{ímpar},\text{par}\}$: semestre da turma;
    \item $\mathcal{P}_c\subseteq\mathcal{P}$: professores habilitados para $c\in\mathcal{C}^{\mathrm{IC}}$;
    \item $\mathcal{Q}_c\subseteq\mathcal{Q}$: padrões admissíveis para $c$;
    \item $\mathcal{R}_{c,m}\subseteq\mathcal{R}$: salas admissíveis para o encontro $m$;
    \item $\mathcal{C}_g\subseteq\mathcal{C}$: turmas pertencentes ao grupo curricular $g$;
    \item $a_{c,m,h,q}\in\{0,1\}$: indica se o encontro $m$ ocupa o slot $h$ quando o padrão $q$ é escolhido.
\end{itemize}

Parâmetros principais:
\begin{itemize}
    \item $n_c$: número de vagas da turma;
    \item $\operatorname{cap}_r$: capacidade da sala;
    \item $\rho_{c,m}$: indicador de exigência de laboratório ou recurso especial;
    \item $\ell_r$: indicador de sala compatível com laboratório ou recurso;
    \item $\operatorname{pref}_{p,c}\in[0,1]$: preferência de $p$ pela disciplina de $c$;
    \item $\pi_p\in[0,1]$: prioridade de $p$;
    \item $\bar q_c$: padrão fixo de $c\in\mathcal{C}^{\mathrm{fix}}$;
    \item $\bar p_c$: professor exógeno de $c\in\mathcal{C}^{\mathrm{out}}$; não é variável do IC;
    \item $\operatorname{desc}_{\min}$: descanso mínimo entre jornadas;
    \item $w_{\mathrm{pref}},w_{\mathrm{dias}},w_{\mathrm{jan}},w_{\mathrm{cap}},w_{\mathrm{rod}}\geq0$: pesos da função objetivo; \pendente{definir os valores dos pesos.}
\end{itemize}

\subsection{Variáveis de decisão}
\label{subsec:variaveis}

Para turmas do IC:
\begin{align}
    x_{c,p} &= \begin{cases}1,&\text{se }c\text{ é atribuída a }p,\\0,&\text{caso contrário,}\end{cases}
    &&c\in\mathcal{C}^{\mathrm{IC}},\ p\in\mathcal{P}_c,\\
    y_{c,q} &= \begin{cases}1,&\text{se }c\text{ recebe o padrão }q,\\0,&\text{caso contrário,}\end{cases}
    &&c\in\mathcal{C},\ q\in\mathcal{Q}_c,\\
    z_{c,m,r} &= \begin{cases}1,&\text{se o encontro }m\text{ de }c\text{ usa }r,\\0,&\text{caso contrário,}\end{cases}
    &&c\in\mathcal{C}^{\mathrm{IC}},\ m\in\mathcal{M}_c,\ r\in\mathcal{R}_{c,m}.
\end{align}

A ocupação da turma e do encontro em um slot é dada por:
\begin{align}
    u_{c,h} &= \sum_{q\in\mathcal{Q}_c:\,h\in H_q} y_{c,q},\\
    v_{c,m,h} &= \sum_{q\in\mathcal{Q}_c} a_{c,m,h,q}y_{c,q}.
\end{align}

Para turmas externas, professor e horário são parâmetros. Essas turmas entram nas verificações curriculares como ocupações fixas, sem criar variáveis de professor externo no conjunto $\mathcal{P}$.

\subsection{Restrições fortes propostas}
\label{subsec:restricoes_hard}

\paragraph{(H1) Professor único.}
\begin{equation}
    \sum_{p\in\mathcal{P}_c}x_{c,p}=1,
    \qquad \forall c\in\mathcal{C}^{\mathrm{IC}}.
\end{equation}

\paragraph{(H2) Padrão único de horário.}
\begin{equation}
    \sum_{q\in\mathcal{Q}_c}y_{c,q}=1,
    \qquad \forall c\in\mathcal{C}.
\end{equation}

\paragraph{(H3) Sala única por encontro.}
\begin{equation}
    \sum_{r\in\mathcal{R}_{c,m}}z_{c,m,r}=1,
    \qquad \forall c\in\mathcal{C}^{\mathrm{IC}},\ m\in\mathcal{M}_c.
\end{equation}

\paragraph{(H4) Horário fixo.}
\begin{equation}
    y_{c,\bar q_c}=1,
    \qquad \forall c\in\mathcal{C}^{\mathrm{fix}}.
\end{equation}
O professor das turmas externas é o parâmetro $\bar p_c$ e não participa de $x$.

\paragraph{(H5) Domínio de setor.} Para uma turma do IC com horário variável, $\mathcal{Q}_c$ contém apenas os padrões compatíveis com os dias autorizados para seu setor. \pendente{preencher os domínios com os dias e habilitações definitivos dos setores.}

\paragraph{(H6) Ausência de choque de professor.}
\begin{equation}
    \sum_{c\in\mathcal{C}^{\mathrm{IC}}:\,\sigma(c)=\sigma}x_{c,p}u_{c,h}\leq1,
    \quad \forall p\in\mathcal{P},\ h\in\mathcal{H},\ \sigma\in\{\text{ímpar},\text{par}\}.
\end{equation}

\paragraph{(H7) Ausência de choque de sala.}
\begin{equation}
    \sum_{c\in\mathcal{C}^{\mathrm{IC}}:\,\sigma(c)=\sigma}
    \sum_{m\in\mathcal{M}_c}z_{c,m,r}v_{c,m,h}\leq1,
    \quad \forall r,h,\sigma.
\end{equation}

\paragraph{(H8) Ausência de choque curricular.}
\begin{equation}
    \sum_{c\in\mathcal{C}_g}u_{c,h}\leq1,
    \qquad \forall g\in\mathcal{G},\ h\in\mathcal{H}.
\end{equation}
Cada $g$ contém somente turmas do mesmo semestre. Turmas externas com horário fixo também participam de $\mathcal{C}_g$ quando pertencem ao período.

\paragraph{(H9) Compatibilidade de recursos.}
\begin{equation}
    z_{c,m,r}=0\quad\text{se }\rho_{c,m}=1\text{ e }\ell_r=0.
\end{equation}
\pendente{definir se encontros sem exigência de laboratório podem ocupar laboratórios ou se essa utilização deve ser proibida ou penalizada.}

\paragraph{(H10) Capacidade.}
\begin{equation}
    z_{c,m,r}=0\quad\text{sempre que }n_c>\operatorname{cap}_r.
\end{equation}

\paragraph{(H11) Descanso entre jornadas.} Para cada semestre $\sigma$ e cada par $(h,h')$ de dias consecutivos que viole o descanso mínimo:
\begin{equation}
    \sum_{c\in\mathcal{C}^{\mathrm{IC}}:\,\sigma(c)=\sigma}x_{c,p}u_{c,h}
    +\sum_{c\in\mathcal{C}^{\mathrm{IC}}:\,\sigma(c)=\sigma}x_{c,p}u_{c,h'}\leq1,
    \quad \forall p\in\mathcal{P}.
\end{equation}

\paragraph{(H12) Carga anual mínima.}
\begin{equation}
    \sum_{c\in\mathcal{C}^{\mathrm{IC}}\cap\mathcal{C}^{\mathrm{obr}}}x_{c,p}\geq3,
    \qquad \forall p\in\mathcal{P}^{\mathrm{H12}}.
\end{equation}
A restrição utiliza o mínimo anual de três obrigatórias. \pendente{definir a composição de $\mathcal{P}^{\mathrm{H12}}$ e registrar a fonte e a unidade de contagem da regra.}

\paragraph{Observação sobre linearização.} As expressões H6, H7 e H11 usam produtos de variáveis binárias para apresentar a lógica de forma compacta. Caso a formulação seja convertida em um modelo linear inteiro, serão necessárias variáveis auxiliares e restrições de ligação. Essa linearização não faz parte do protótipo metaheurístico atual.

\subsection{Função objetivo proposta}
\label{subsec:funcao_objetivo}

A função objetivo conceitual é:
\begin{equation}
\min Z=-w_{\mathrm{pref}}\Phi_{\mathrm{pref}}
+w_{\mathrm{dias}}\Phi_{\mathrm{dias}}
+w_{\mathrm{jan}}\Phi_{\mathrm{jan}}
+w_{\mathrm{cap}}\Phi_{\mathrm{cap}}
+w_{\mathrm{rod}}\Phi_{\mathrm{rod}}.
\end{equation}

Seus componentes são:
\begin{itemize}
    \item $\Phi_{\mathrm{pref}}=\sum_{c,p}\pi_p\operatorname{pref}_{p,c}x_{c,p}$: atendimento de preferências ponderado pela prioridade;
    \item $\Phi_{\mathrm{dias}}$: número de combinações professor--semestre--dia com pelo menos uma aula;
    \item $\Phi_{\mathrm{jan}}$: quantidade de slots ociosos entre aulas do mesmo professor no mesmo dia;
    \item $\Phi_{\mathrm{cap}}=\sum_{c,m,r}z_{c,m,r}(\operatorname{cap}_r-n_c)$: desperdício de capacidade, quando H10 é satisfeita;
    \item $\Phi_{\mathrm{rod}}$: repetição do professor na mesma disciplina entre os dois semestres do ano. Uma comparação com anos anteriores poderá ser acrescentada quando os dados e a regra de rodízio forem definidos.
\end{itemize}

Na configuração inicial de implementação, os critérios suaves são somados com peso unitário e as violações fortes recebem penalidade de $10^6$. \pendente{definir a classificação final de H8, H10, H11 e H12 e substituir os pesos de validação pelos valores usados nos experimentos.}

\section{Proposta de Solução e Implementação}
\label{sec:proposta_solucao}

A solução é organizada em componentes de representação, avaliação, vizinhança e busca. Os componentes possuem implementações iniciais em Python, usadas para verificar seu comportamento de forma isolada. \pendente{concluir o registro desses componentes no pyoptframe e atualizar esta seção com a arquitetura integrada.}

\subsection{Representação e organização}
\label{subsec:arquitetura}

A instância é armazenada em JSON. Cada turma contém, entre outros campos, o professor atual, os encontros semanais, as salas, a classificação curricular e, quando disponíveis, os domínios de professores e horários. Durante a avaliação, esses dados são convertidos em tabelas de turmas e encontros.

O código está dividido principalmente em:
\begin{itemize}
    \item \texttt{src/eval/}: leitura da instância, inferência de recursos, tratamento de salas, preferências e avaliador de referência;
    \item \texttt{src/solve/}: avaliador leve e implementações manuais de SA para salas, salas e horários, e professores.
\end{itemize}

Horários marcados como fixos não pertencem à lista de turmas flexíveis do solver de horários. Os domínios usados na implementação inicial foram derivados do histórico de 2025. \pendente{substituir os domínios inferidos pelos dados definitivos de setores e habilitações.}

\subsection{Avaliadores}
\label{subsec:avaliador}

O avaliador de referência, implementado em \texttt{src/eval/evaluator.py}, usa pandas e produz uma decomposição por critério. Ele contabiliza atualmente:
\begin{itemize}
    \item conflitos de sala, professor e grupo curricular;
    \item capacidade insuficiente e incompatibilidade de laboratório;
    \item descanso insuficiente e professores abaixo do mínimo H12;
    \item dias trabalhados, janelas, desperdício de capacidade, rodízio e preferência ponderada.
\end{itemize}

As condições estruturais H1--H5 são tratadas principalmente pela construção da instância e pelos domínios dos movimentos. \pendente{adicionar contadores independentes para as condições estruturais ainda não cobertas, incluindo encontros sem sala em H3.}

Para reduzir o custo causado pelas operações de pandas durante a busca, \path{src/solve/direct_objective.py} implementa a mesma avaliação diretamente sobre o JSON. Esse avaliador é mais leve, mas ainda recalcula a solução completa a cada movimento; não se trata de uma avaliação delta. Testes automatizados verificam a equivalência entre os dois avaliadores para os casos cobertos.

O score usado na configuração inicial é:
\begin{equation}
    F(s)=10^6\,V_{\mathrm{hard}}(s)+\sum_j\Phi_j(s),
\end{equation}
onde $V_{\mathrm{hard}}$ é a soma dos sete contadores fortes implementados e os critérios suaves disponíveis entram com peso unitário. Essa escala serve para testar a prioridade das violações fortes, mas não corresponde aos pesos finais da modelagem.

\subsection{Movimentos implementados}
\label{subsec:vizinhancas}

Foram implementados três tipos de alteração:
\begin{enumerate}
    \item \textbf{sala}: altera a sala de um encontro para outra candidata compatível com a capacidade disponível e a classificação de laboratório;
    \item \textbf{horário}: substitui os encontros de uma turma flexível por outro padrão pertencente ao domínio histórico da turma;
    \item \textbf{professor}: realoca uma turma para outro nome presente no domínio histórico de professores observado para seu setor.
\end{enumerate}

\pendente{implementar a troca simultânea de professores, a troca de salas entre turmas e a composição dos três tipos de decisão em uma mesma busca.}

\subsection{Implementação inicial de Simulated Annealing}
\label{subsec:sa_solver}

Os solvers atuais implementam o SA diretamente em Python. Eles partem da alocação histórica contida na instância, usam um número fixo de iterações e reduzem linearmente a temperatura:
\begin{equation}
    T(i)=\max\left\{1,T_0\left(1-\frac{i}{N}\right)\right\},
\end{equation}
onde $N$ é o total de iterações. O solver conjunto de salas e horários escolhe um movimento de horário com probabilidade $0{,}6$ e um movimento de sala com probabilidade $0{,}4$. O solver de professores executa apenas realocações individuais.

O fluxo comum é apresentado no Algoritmo~\ref{alg:sa_atual}.

\begin{algorithm}[H]
\SetAlgoLined
\KwIn{Instância $I$, temperatura $T_0$, iterações $N$ e semente}
\KwOut{Melhor solução encontrada segundo $F$}
$s\gets\text{cópia da alocação inicial de }I$\;
$s^*\gets s$\;
\For{$i\gets0$ \textbf{até} $N-1$}{
    sortear um movimento permitido e aplicá-lo provisoriamente\;
    $\Delta\gets F(s')-F(s)$\;
    $T\gets\max\{1,T_0(1-i/N)\}$\;
    \eIf{$\Delta\leq0$ \textbf{ou} $\operatorname{Random}(0,1)<\exp(-\Delta/T)$}{
        $s\gets s'$ e atualizar $s^*$ se necessário\;
    }{
        desfazer o movimento\;
    }
}
\Return{$s^*$}\;
\caption{Estrutura da implementação inicial de SA}
\label{alg:sa_atual}
\end{algorithm}

O algoritmo retorna a melhor solução segundo o score, que pode continuar inviável quando nenhuma solução viável é alcançada durante a execução.

\subsection{Integração com pyoptframe}

A implementação no \textit{pyoptframe} utiliza um construtivo, um avaliador e as estruturas de vizinhança correspondentes como componentes registrados no motor do framework. \pendente{implementar a integração, indicar as classes e identificadores empregados e descrever a configuração de SA, ILS e/ou VNS efetivamente comparada.}

\section{Experimentos Computacionais}
\label{sec:experimentos}

Esta seção apresenta as instâncias, o protocolo de execução, os cenários de comparação e os resultados computacionais.

\subsection{Instâncias}
\label{subsec:instancias}

\pendente{descrever as instâncias definitivas: origem, quantidade de turmas, professores, salas, encontros, grupos curriculares e parâmetros de cada uma.}

A pipeline produziu uma instância de validação a partir do quadro de horários de 2025. O arquivo \texttt{instancia\_2025\_cc\_si.json} contém 154 turmas, das quais 150 foram classificadas como turmas do IC, 58 nomes no universo inicial de professores e 20 salas. Essa instância foi usada para verificar a leitura dos dados e o comportamento dos componentes do solver.

Na instância de validação, capacidades de sala e exigências de laboratório foram inferidas do histórico, as prioridades foram fixadas em $1{,}0$ e os domínios de professor e horário foram derivados dos setores observados em 2025. Esses parâmetros não devem ser confundidos com dados institucionais definitivos.

\pendente{substituir capacidades, recursos, prioridades, setores, habilitações e universo de H12 pelos valores utilizados nos experimentos finais.}

A configuração contém 148 turmas classificadas como obrigatórias. Se os 58 professores forem incluídos em $\mathcal{P}^{\mathrm{H12}}$, a restrição exigirá pelo menos $3\times58=174$ atribuições de obrigatórias. Esse universo, portanto, torna a instância de validação estruturalmente inviável para H12.

\subsection{Protocolo experimental}
\label{subsec:protocolo}

Para cada combinação de instância e algoritmo, serão realizadas execuções independentes com sementes distintas. As soluções serão comparadas pela quantidade de violações fortes, pelo valor da função objetivo e pela decomposição dos critérios suaves.

\begin{itemize}
    \item \textbf{Número de execuções:} \pendente{definir.}
    \item \textbf{Critério de parada ou limite de tempo:} \pendente{definir.}
    \item \textbf{Sementes:} \pendente{definir ou indicar o procedimento de geração.}
    \item \textbf{Máquina e ambiente de execução:} \pendente{informar processador, memória, sistema, Python e versões das bibliotecas.}
    \item \textbf{Parâmetros do SA:} \pendente{informar temperatura inicial, esquema de resfriamento e número de iterações.}
    \item \textbf{Parâmetros de ILS/VNS:} \pendente{preencher para os métodos efetivamente utilizados.}
\end{itemize}

O baseline histórico será construído a partir do quadro observado e avaliado pelos mesmos critérios aplicados às soluções geradas. \pendente{descrever a auditoria do baseline e o tratamento de salas ausentes, turmas paralelas e registros duplicados.}

\subsection{Cenários de integração entre CC e SI}

O acoplamento entre as grades será analisado em três cenários:
\begin{itemize}
    \item \textbf{E1 (CC $\rightarrow$ SI)}: otimizar CC, congelar suas decisões e, em seguida, otimizar as decisões ainda livres de SI;
    \item \textbf{E2 (SI $\rightarrow$ CC)}: executar a ordem inversa;
    \item \textbf{E3 (conjunto)}: otimizar simultaneamente as decisões de ambas as grades.
\end{itemize}

\pendente{definir o tratamento das turmas compartilhadas em E1 e E2 e as métricas usadas para medir o custo da ordem de planejamento.}

\subsection{Validação inicial da implementação}
\label{subsec:validacao_inicial}

Antes do protocolo comparativo, foram realizadas execuções determinísticas com semente 2025 para verificar o efeito dos movimentos sobre o avaliador. Esses testes têm função de validação de software e não integram a comparação final entre algoritmos.

\begin{table}[htbp]
\centering
\caption{Validação inicial na instância CC/SI 2025}
\label{tab:resultados_preliminares}
\small
\setlength{\tabcolsep}{3pt}
\begin{tabular}{lrrrrrrr}
\toprule
\textbf{Configuração} & \textbf{Score} & \textbf{Hard} & \textbf{Sala} & \textbf{Curr.} & \textbf{H12} & \textbf{Janelas} & \textbf{Desp. cap.} \\
\midrule
Baseline extraído & 51.002.369 & 51 & 13 & 10 & 28 & 28 & 2.172 \\
SA -- salas & 38.004.258 & 38 & 0 & 10 & 28 & 28 & 4.060 \\
SA -- salas e horários & 29.004.412 & 29 & 0 & 1 & 28 & 80 & 4.144 \\
SA -- professores & 35.002.435 & 35 & 13 & 10 & 12 & 59 & 2.172 \\
\bottomrule
\end{tabular}
\end{table}

O movimento de sala eliminou os conflitos de sala contabilizados na instância de validação sem alterar H12. O movimento de horário reduziu os conflitos curriculares de 10 para 1, mas aumentou janelas e desperdício estimado. O movimento de professor reduziu de 28 para 12 a quantidade de professores abaixo do mínimo, sem alcançar viabilidade para H12. Esses resultados confirmam quais parcelas do avaliador são afetadas por cada movimento, mas não estabelecem superioridade sobre o baseline.

\subsection{Resultados comparativos}

\pendente{inserir tabelas e gráficos das execuções finais, medidas agregadas, análise estatística, comparação com o baseline e análise dos cenários E1--E3.}

\section{Conclusão}
\label{sec:conclusao}

\pendente{redigir após os experimentos: retomar o objetivo geral, sintetizar a modelagem e a implementação, apresentar os principais resultados e indicar as limitações do estudo.}

\subsection{Trabalhos futuros}

\pendente{definir a partir das limitações observadas nos experimentos. Possibilidades: novas vizinhanças, métodos híbridos, calibração automática de parâmetros, atualização dos dados e interface de apoio à visualização das soluções.}


% --- REFERÊNCIAS BIBLIOGRÁFICAS ---

\newpage
\bibliographystyle{plain}
\bibliography{../referencias/referencias}

\end{document}
